El \textit{Meteosat} és un satèl·lit meteorològic llançat per l'Agència Espacial Europea (ESA) que proporciona informació meteorològica d'Àfrica i Europa. Com que l'objectiu del \textit{Meteosat} és oferir imatges d'una mateixa zona del planeta, el satèl·lit segueix una òrbita geostacionària: gira en el pla equatorial a la mateixa velocitat angular que la Terra.

\figura[0.45]{fig1.png}

\begin{apartat}{1}
A quina distància de la superfície terrestre es troba el \textit{Meteosat}?
\end{apartat}

\begin{apartat}{1}
Quina és l'energia cinètica del \textit{Meteosat}? Quina energia mínima caldria proporcionar-li perquè s'allunyés indefinidament de la Terra?
\end{apartat}

\dades{%
  $G = 6,67\times 10^{-11}\ \mathrm{N\,m^2\,kg^{-2}}$\\
  $R_\mathrm{Terra} = 6\,370\ \mathrm{km}$\\
  $M_\mathrm{Terra} = 5,97\times 10^{24}\ \mathrm{kg}$\\
  $m_\textit{Meteosat} = 2,00\times 10^{3}\ \mathrm{kg}$}
